% TOP_PROBLEM: {"id": 1251, "title": "Almost Perfect Numbers Beyond Powers of 2", "category_id": 1, "category": {"id": 1, "name": "number_theory", "display_name": "Number Theory", "description": "Properties of integers, prime numbers, Diophantine equations.", "slug": "number-theory", "order_index": 1, "created_at": "2026-07-31T15:26:25.670Z"}, "status": "open"}
% ATTEMPT_STATUS: unresolved
% Research attempt by GPT_6_Astra_Ultra, 1 October 2026.
% Canonical UnsolvedMath ID; original statements and sources preserved.
% FIRST_POSTED: 2026-10-01
% Imported from: unsolvedmath_additional_500.tex; UnsolvedMath ID 1251
% !TeX program = lualatex
% Compile twice: lualatex 1251.tex
% Self-contained: no external bibliography, images, or \input files.
% Numeric section labels and visible numbers are original UnsolvedMath IDs.
\documentclass[11pt,a4paper]{article}
\usepackage[margin=24mm,headheight=14pt]{geometry}
\usepackage{fontspec}
\setmainfont{Latin Modern Roman}
\setsansfont{Latin Modern Sans}
\setmonofont{Latin Modern Mono}
\usepackage{amsmath,amssymb,mathtools}
\usepackage{unicode-math}
\setmathfont{Latin Modern Math}
\usepackage{microtype}
\usepackage{enumitem,xurl,needspace}
\usepackage[dvipsnames]{xcolor}
\usepackage{hyperref}
\hypersetup{unicode=true,colorlinks=true,linkcolor=MidnightBlue,urlcolor=MidnightBlue,
 pdftitle={UnsolvedMath Problem 1251},pdfauthor={GPT 6 Astra Ultra},bookmarksnumbered=true}
\usepackage{bookmark,tocloft,titlesec,fancyhdr}
\setlength{\cftsecnumwidth}{4.2em}
\cftsetpnumwidth{2.5em}
\cftsetrmarg{4em}
\titleformat{\section}{\Large\bfseries\raggedright}{\thesection}{0.7em}{}
\pagestyle{fancy}\fancyhf{}
\fancyhead[L]{\small\sffamily GPT 6 Astra Ultra research attempt}
\fancyhead[R]{\small Original catalogue IDs}
\fancyfoot[C]{\thepage}
\setlength{\parindent}{0pt}
\setlength{\parskip}{5pt plus 1pt}
\setlength{\emergencystretch}{3em}
\setcounter{tocdepth}{1}\setcounter{secnumdepth}{1}
\setlist[itemize]{leftmargin=3em,itemsep=4pt,topsep=4pt,parsep=0pt}
\allowdisplaybreaks
\newcommand{\N}{\mathbb{N}}
\newcommand{\Z}{\mathbb{Z}}
\newcommand{\Q}{\mathbb{Q}}
\newcommand{\R}{\mathbb{R}}
\newcommand{\C}{\mathbb{C}}
\newcommand{\F}{\mathbb{F}}
\newcommand{\A}{\mathbb{A}}
\newcommand{\PP}{\mathbb{P}}
\newcommand{\E}{\mathbb{E}}
\newcommand{\eps}{\varepsilon}
\newcommand{\dd}{\,\mathrm d}
\newcommand{\sref}[2]{\hyperlink{source:#1:#2}{[S#2]}}
\newif\ifshowcatalogue
\showcataloguetrue
\newcommand{\cataloguescope}[1]{\ifshowcatalogue
 \paragraph{Statement recorded in the supplied source.}
 #1\fi}
\begin{document}
% UnsolvedMath ID 1251; source catalogue code NT-044

\setcounter{section}{1250}

\section{Almost Perfect Numbers Other Than Powers of Two}\label{1251}

{\small\sffamily UnsolvedMath ID 1251 · Number Theory}

\subsection{Definitions and mathematical statement}

\paragraph{Shared notation.}Unless a section says otherwise, $\N=\{1,2,\ldots\}$,
$\Z$ is the ring of integers, $\Q,\R,\C$ are the rational, real, and complex
fields, $[N]=\{1,\ldots,\lfloor N\rfloor\}$, and $\log$ is the natural
logarithm. For real functions of a parameter tending to infinity,
$f\ll g$ and $f=O(g)$ mean $|f|\leq Cg$ eventually for some fixed $C$;
$f\gg g$ reverses this comparison; $f=o(g)$ means $f/g\to0$;
$f\sim g$ means $f/g\to1$; and $f\asymp g$ means both order bounds.
A subscript on $O$ or $\ll$ specifies allowed constant dependence.
The expression $n^{o(1)}$ in an upper bound means growth smaller than
$n^\epsilon$ for each fixed $\epsilon>0$. Local definitions govern any
exception, finite-field convention, or use of the same symbol for another
object. An asymptotic claim is not automatically an effective algorithm.



An almost perfect number is a positive integer $n$ with $\sigma(n)=2n-1$. The powers $2^k$, including $2^0=1$, satisfy this equality. The target is the existence or impossibility of a solution outside $\{2^k:k\geq0\}$. \sref{1251}{1} \sref{1251}{2}

\paragraph{Statement recorded in the supplied source.}
Do any almost perfect numbers exist that are not powers of 2? \sref{1251}{1}

\paragraph{Residual task reported by the archive.}

The embedded review (2026-08-17) records: Construct an almost perfect number outside the powers of two or exclude all remaining structural cases. \sref{1251}{1}

\subsection{Short English statement}

Apart from powers of two, can any positive integer have proper divisors whose sum falls short of that integer by exactly one?

\subsection{Research attempt}

\paragraph{Scope and approach.}
The equation is $\sigma(n)=2n-1$, so the powers of two, including one,
are already solutions. The question is whether there are any others.
The Antalan--Dris primary preprint [S2], version 4, was inspected at
abstract level on 1 October 2026; its stated result concerns structural
restrictions on a hypothetical even example, not a construction.
This attempt derives a cofactor identity, eliminates every number with
just one distinct odd prime factor, and isolates the remaining problem.

\paragraph{The parity reduction and all known immediate examples.}
The right side is odd. For an odd prime $p$, the divisor sum
$1+p+\cdots+p^e$ is odd exactly when $e$ is even. Therefore
\[
                          n=2^a m^2
\]
with $a\ge0$ and $m$ odd. Conversely this shape merely ensures odd
divisor sum; it does not imply the required equality. If $m=1$, then
$\sigma(2^a)=2^{a+1}-1=2n-1$, proving all the stated power-of-two
examples, including $a=0$. We henceforth take $m>1$.

\paragraph{The exact odd-cofactor identity.}
Let $A=2^{a+1}$ and $D(m^2)=\sigma(m^2)-m^2$, the sum of the
proper positive divisors of $m^2$. Multiplicativity gives
\[
 (A-1)\sigma(m^2)=Am^2-1
 \quad\Longleftrightarrow\quad
                       (A-1)D(m^2)=m^2-1.
\]
Thus $A-1$ divides $m^2-1$, and $\gcd(m,A-1)=1$.
For a prescribed odd $m>1$, the only possible power-of-two part is
determined by
\[
             A=1+\frac{m^2-1}{D(m^2)}.
\]
The quotient must be an integer and the resulting $A$ must be a power
of two at least two. Conversely, if these tests pass, the original
equation follows, so this is a complete criterion for each cofactor.
There is at most one exponent $a$ for a fixed odd $m>1$.

The sign matters: the analogous quasiperfect equation would have
$m^2+1$ on the right, producing a sum-of-two-squares obstruction.
Here $m^2-1=(m-1)(m+1)$ has no such obstruction. Importing that
plus-sign argument into this problem would incorrectly eliminate the
known powers of two and the unresolved cases.

\paragraph{A complete infinite exclusion: one odd prime factor.}
Suppose $m=p^e$ with $p$ odd prime and $e\ge1$. The proper-divisor
sum is
\[
 D(p^{2e})=1+p+\cdots+p^{2e-1}
                  =\frac{p^{2e}-1}{p-1}.
\]
Substitution into the cofactor identity and cancellation of the positive
$p^{2e}-1$ gives $A-1=p-1$, or $p=A=2^{a+1}$. No odd prime is a
power of two. Thus no solution outside the powers of two can have only
one distinct odd prime factor. This includes all odd prime powers and
all integers of the form $2^a p^{2e}$, with no bound on the exponents.
Every nontrivial candidate needs at least two distinct odd primes.

\paragraph{A useful size restriction with proof.}
Since $m>1$, both $1$ and $m$ are distinct proper divisors of $m^2$.
Hence $D(m^2)\ge m+1$. The cofactor identity then implies
\[
               A-1\le\frac{m^2-1}{m+1}=m-1,
 \qquad A\le m.
\]
But $A$ is even and $m$ is odd, so equality is impossible and
\[
                         2^{a+1}<m.
\]
This recovers directly the kind of structural inequality stated in
[S2], without using its proof. It bounds the even part in terms of the
odd square root, not the total size of the candidate. In particular it
does not imply that only finitely many $m$ can occur.

There is a residue restriction on factors of $m$ as well. If $3\mid m$,
then $3\nmid A-1$ by coprimality, so $a+1$ cannot be even, because
$2^{a+1}\equiv1\pmod3$ for even $a+1$. Thus $a$ must be even on
that branch. More generally the order of two modulo any odd prime
dividing $m$ cannot divide $a+1$. These restrictions discard some
exponents but do not exhaust all exponents simultaneously.

\paragraph{Abundancy and the limiting-product issue.}
Dividing the defining equation by $n$ gives
\[
 \frac{\sigma(n)}n=2-\frac1n<2,
 \qquad
 \frac{\sigma(m^2)}{m^2}
       =\frac{A-m^{-2}}{A-1}.
\]
For fixed $A$, the latter is just below $A/(A-1)$, with a prescribed
error $1/((A-1)m^2)$. Having a deficient integer, meaning divisor sum
less than twice itself, is much weaker than attaining deficiency
exactly one. A sequence of abundancy values approaching two from below
does not prove that any one of them has the required exact error.

For a fixed prime support, upper products such as
$\prod p/(p-1)$ help bound possible exponents, but the support is not
fixed in the original problem. An argument that exhausts a finite list
of supports needs a proof that every candidate lies on that list.
Neither the square-part reduction nor $A<m$ provides such a list.

\paragraph{Executed finite classification.}
The script evaluates $\sigma(n)$ exactly for $1\le n\le200000$.
The solutions to $\sigma(n)=2n-1$ are exactly
\[
                    1,2,4,\ldots,131072,
\]
namely the powers $2^a$ with $0\le a\le17$. Each integer in the
range is tested, rather than sampling only odd squares or a restricted
exponent pattern. The script and result are \texttt{checks.py} and
\texttt{results.json} in
\path{_Local/Erdos_problems/UnsolvedMath/attempt_audit_2026_10_01/number_theory/}.
The finite classification is independent of the proof of the
one-odd-prime exclusion, which covers an infinite family.

\paragraph{Remaining gap and outcome.}
For a nontrivial candidate the odd part must be a square with at least
two distinct prime factors, its cofactor quotient must give exactly a
power of two, and $2^{a+1}<m$. No argument here rules out all such
cofactors or produces one that passes. \textbf{Outcome: UNRESOLVED.}
The work gives exact reductions and an infinite excluded subclass,
without asserting that the finite list of powers of two is globally
complete.

\subsection{Sources}

{\small

\begin{itemize}

\item[\hypertarget{source:1251:1}{S1}] ULAM.AI, \emph{UnsolvedMath}, supplied \texttt{problems.json}, record 1251 (NT-044), statement and background. The embedded literature-review date is 2026-08-17; that is the archive’s review date, not a fresh certification. Dataset landing page: \url{https://huggingface.co/datasets/ulamai/UnsolvedMath}.

\item[\hypertarget{source:1251:2}{S2}] John Rafael M. Antalan and Jose Arnaldo B. Dris, Some New Results On Even Almost Perfect Numbers Which Are Not Powers Of Two, arXiv:1602.04248. \url{https://arxiv.org/abs/1602.04248}. Bibliographic information imported from the supplied record.

\item[\hypertarget{source:1251:3}{S3}] Encyclopedia of Mathematics, Almost perfect number. \url{https://encyclopediaofmath.org/wiki/Almost_perfect_number}. Bibliographic information imported from the supplied record.


\item[E1] John Rafael M. Antalan and Jose Arnaldo B. Dris,
\emph{Some New Results On Even Almost Perfect Numbers Which Are Not
Powers Of Two}, arXiv:1602.04248v4 (2017).
\url{https://arxiv.org/abs/1602.04248v4}. Abstract inspected
1 October 2026. The inequality $2^{a+1}<m$ used above is proved
directly in this notebook as well as reported in that source.

\end{itemize}

}

\label{end:1251}
\end{document}
